\chapter{Implementation}
\label{ch:implementation}


This chapter is going to focus on the details of the process of creating the playable game
that used voice as a potential input modality. The first step in the process was to identify the type of tools
that were going to be needed in the game. 
First of all there needed to be the video game itself. 

% 2D and 3D
The most fundemental choice for the game prototype is if it should use 2 dimentional or 3D graphics. This 
choice might seem trivial but the consequences are very far fetching. 

% complexity of implementation
Adding the third dimention in game envirornment is adding additional layer of complexity. Instead of dealing with $x$ and $y$ 
there's also $z$ creating a 3 dimentional system. That in itself makes it that the assets in the game have to be also three dimentional.
Tha in itself means usually longer time needed to create projects since the process of creating assets becomes longer and more advanced.
Another thing that makes making 3D games more difficult is physics. The vast majority of games have some sort of physics 
implementation, be it collision detection, particle systems or force vectors. Making those calculations is something that becomes much harder when adding
a third dimension. 
To be able to verify this thesis a quick iteration and implementation are key, being arguments for choosing 2D.



% precision
for this project. When it comes to input in games, one of the most important factors is the precision of the controls - 
since the player actions should align with the expectations of the player.
Given that precision of input is such an important factor the 2D system seems to be making more sense in this scenario,
since it allows to precisely describe things such as movement or rotation using natural language. 
The difference between 2D static camera and 3D first person movement commands for example:

2D: Rotate to the right and move 10 steps forward

3D: Look to the upper left a bit and then move 10 stept forward

This example illustrates that it might be difficult for the player to describe camera motions since 3 dimentional transforms 
are not intuitive and might require repetetive attempts

% why 2D
Of course not all of 3D games rely on a player controllable camera. Many games employ a fixed camera or don't let the player
control the camera at all, but given the time frame of this thesis and the fact that it includes exploring different potential 
prototypes I made the decision that 2D is the preferable option.
After making that choice the next logical step is to think of potential suitable game ideas that would differ from each other
and also would make it possible to utilise the voice input.


% generes a
There's a lot of different video game generes. I set out to find out what type of game would be the most suitable


% Engines
Games are usually built using some sort of game engine, 
allowing for standing on the shoulders of giants so to speak. Engines vary greatly in their complexity, supported platforms,
pricing, licensing, and many other technical factors. For this work I decided to  


\section{WhisperCPP}

How to download custom whispercpp models
\url{https://github.com/ggml-org/whisper.cpp/blob/master/models/README.md}


$large-v3-turbo-q5_0$

The trick to remove silence in ffmpeg

\begin{lstlisting}[style=myStyle, language=bash]
ffmpeg -i fox.mp4 -ar 16000 -ac 1 -c:a pcm_s16le -af silenceremove=1:0:-50dB fox.wav
\end{lstlisting}



whisper cpp flag - can be 32 or 0 too but quality suffers
\begin{lstlisting}[style=myStyle, language=bash]
    --max-context 64 --entropy-thold 2.8
\end{lstlisting}


Running Whispercpp on my CPU proved to be remarkably slow


\url{https://old.reddit.com/r/LocalLLaMA/comments/1fyvc60/how_to_improve_whisper_translation_it_keeps/}


Discussion about whisper cpp 
\url{https://old.reddit.com/r/LocalLLaMA/comments/1hc1qzi/is_whispercpp_still_the_king_of_stt/}

whisper cpp wav file to txt file example
\begin{lstlisting}[style=myStyle, language=bash]
./build/bin/whisper-cli -m models/ggml-large-v3-turbo-q5_0.bin -nt samples/jfk.wav > test.txt
\end{lstlisting}

whisper cpp server run command
\begin{lstlisting}[style=myStyle, language=bash]

./build/bin/whisper-server -m models/ggml-large-v3-turbo-q5_0.bin -nt --host 0.0.0.0 -debug --port 8008 -t 7
\end{lstlisting}

Document the process of building whisper cpp 
There's a screenshot in the notes folder
\section{LLAMACPP}

Running LLAma CPP server
\begin{lstlisting}[style=myStyle, language=bash]
    ./llama-server -m ../../models/ToolACE-2-Llama-3.1-8B-Q4_K_M.ggml --n-gpu-layers 60 --host 0.0.0.0 --jinja
\end{lstlisting}

The docs \url{https://github.com/ggml-org/llama.cpp/blob/master/tools/server/README.md}

The models themselves need to be quantified \url{https://github.com/ggml-org/llama.cpp/blob/master/tools/quantize/README.md}

There's a note called Quanting the model with images


Theoretical navigation
Action + object = navigation
Free flowing conversation = talking to NPC

Though, to be said this would require some more testing into what's actually feasible here, but it could be a good starting point, especially the first system.

I made a smimple bash script to automate the deployment of the AI backend when run.

\section{Tool Calling}

Best model that fits on my machine
\url{https://huggingface.co/Team-ACE/ToolACE-2-Llama-3.1-8B
}

There's a nice benchmark for tool calling
\url{https://gorilla.cs.berkeley.edu/leaderboard.html}

\section{Game prototype}

Multiple engines were considered
\begin{itemize}
    \item Unity3D
    
    This engine while very powerful has a huge moving mass of all of it's complexity. This means that while 
    it allows for creating of very advanced titles, for the ideating stage it doesn't allow for as fast prototype
    creation as other choices

    \item Godot
    
    Godot is a universally praised engine with a cult-following. The main reason for that is it's licensing giving 
    the studios a stable pillar to stand on compared to it's biggest direct competitor being Unity3D. 
    Another benetfit that it has is the fact that it seems to be preferable by some when it comes to 2D game development.
    However, it still suffers from how complex it is, making creating prototypes and ideating more difficult and ultimately
    longer than the engines listed below. 

    \item Pico8
    
    The complete anti-thesis of the above choices. This engine allows for the quickest game creation, offering 
    built in tools, uisng very simple lua language and through creative limitations such as 128x128 resolution forcing
    the game developing process to be simple and streamlined. It's biggest weakness is that it doesn't allow for simple 
    http requests rendering it unusable for this work despite the author's personal preference. 
    
    \item LÖVE
    
    Another simple 2D lua engine, this time with actual support for http requests. It handles higher resolution 
    and makes it possible to quickly create prototypes or even fully fledged games such as Balatro. Hovever due to 
    how niche it is it doesn't have nearly as big of a community making support and learning resources scarce in comparasion
    to bigger engines.

    \item pure CLI
    
    If the game is set to be a text based adventure then maybe an engine isn't needed at all? That would potentially allow
    for huge flexibility when it comes to actual implementation but at a cost of no guidance and best practices. Games are
    rarely written to be exprienced via terminal these days, and even if being text based they still render their text via 
    some kind of graphical front end. Another point speaking against using CLI would be the inablity to use sound or other 
    modalities that are not purely text based. On top of that terminal windows often change their size and are hardly standarised
    which could prove creating UI elements difficult. 
    
    \item PyGame
    
    The choice against this engine comes down mostly due to lacking support and how it seemed to have been replaced by 
    Pico8 and LÖVE in the quick game development community, especialyl looking at their release notes and online activity 
    of users. It offers more capabilities than those engines at the cost of added complexity. It sits somewhere between Pico8
    and Unity3D and for those reasons it wasn't selected for this thesis. 


    \item p5js
    
    The final selection. This engine has the advantages of Pico-8 and LÖVE, allowing for lightning fast development while
    also adding suport for high resolution and http requests. This compbined with big community support and many learning
    resources + with how it runs on Javascript sealed it as the ultimate choice for this thesis.
\end{itemize}